% Generated draft: requires mathematical and source review.
\documentclass[11pt]{article}
\usepackage[margin=25mm]{geometry}
\usepackage{fontspec}
\setmainfont{Latin Modern Roman}
\usepackage{amsmath,amssymb,mathtools,unicode-math}
\setmathfont{Latin Modern Math}
\usepackage{enumitem,needspace,xurl,hyperref}
\setlength{\parindent}{0pt}
\setlength{\parskip}{5pt}
\setlength{\emergencystretch}{3em}
\newcommand{\N}{\mathbb{N}}
\newcommand{\Z}{\mathbb{Z}}
\newcommand{\Q}{\mathbb{Q}}
\newcommand{\R}{\mathbb{R}}
\newcommand{\C}{\mathbb{C}}
\newcommand{\F}{\mathbb{F}}
\newcommand{\A}{\mathbb{A}}
\newcommand{\PP}{\mathbb{P}}
\newcommand{\E}{\mathbb{E}}
\newcommand{\eps}{\varepsilon}
\newcommand{\dd}{\,\mathrm{d}}
\begin{document}
\section*{UnsolvedMath ID: 30003562}
\textbf{Source problem code:} catalogue entry
\subsection{Definitions and mathematical statement}
Let \(\widetilde L\subset\R^2\times(-1,1)\) denote a biperiodic alternating link: its diagram is invariant under a rank-two lattice \(\Gamma\) of planar translations, and overcrossings and undercrossings alternate along each strand. Put \(T^2=\R^2/\Gamma\), and denote the toroidal quotient link by
\[
K=\widetilde L/\Gamma\subset T^2\times(-1,1).
\]
The source calls the biperiodic link \(L\) and its toroidally alternating quotient \(\Lambda\); thus \(K\) denotes the source's \(\Lambda\). The additional restriction ``semi-regular'' is retained as a source-specific qualification: the supplied record gives no definition of this class.

For a nonzero Laurent polynomial \(p\in\C[z^{\pm1},w^{\pm1}]\), its logarithmic Mahler measure is
\[
m(p)=\frac{1}{(2\pi)^2}\int_0^{2\pi}\int_0^{2\pi}
\log\left|p(e^{i\theta},e^{i\phi})\right|\,\dd\theta\,\dd\phi.
\]
Here \(p(z,w)\) is intended to be the characteristic polynomial of the toroidal dimer model associated with the link diagram. The supplied record does not specify the construction of that model, its weights, or the normalization of its characteristic polynomial.

\textbf{Conjectural statement, subject to the qualifications below.} For every semi-regular biperiodic alternating link \(\widetilde L\) and its toroidal quotient \(K\), is
\[
\operatorname{vol}\!\left((T^2\times(-1,1))\setminus K\right)
\leq 2\pi\,m(p(z,w))?
\]
Here the volume is intended as the volume of a complete finite-volume hyperbolic metric on the quotient complement, when such a metric exists.

\textbf{Unresolved formulation issues.} The supplied formula instead writes
\[
\operatorname{vol}\!\left((T^2\times(-1,1))\setminus L\right)
\leq 2\pi\,m(p(z,w)),
\]
although it previously uses \(L\) for the biperiodic link and \(\Lambda\) for the quotient. The formulation with \(K\) above explicitly records the natural quotient interpretation; the supplied evidence does not establish that this is the intended identification. It also does not state the precise semi-regularity conditions, the hypotheses guaranteeing the indicated hyperbolic volume, or the polynomial normalization. In particular, replacing \(p\) by \(cp\), for \(c\ne0\), changes \(m(p)\) by \(\log|c|\), so an arbitrary scalar normalization cannot be ignored. These missing specifications prevent a fully self-contained unconditional formulation.
\subsection{Short English statement}
The proposed inequality bounds the hyperbolic volume of a toroidal quotient of a semi-regular biperiodic alternating link by \(2\pi\) times the logarithmic Mahler measure of its associated dimer polynomial. The quotient notation, precise link class, volume hypotheses, and polynomial normalization require clarification. The record labels the conjecture open; this is a supplied status assessment.
\subsection{Sources}
\begin{itemize}
\item Supplied local UnsolvedMath registry record, permanent ID \(30003562\), title ``Volume Bounds from Mahler Measures of Biperiodic Links,'' source problem code \texttt{catalogue entry}, in the Oberwolfach Reports Open Problems collection. The record lists proposed year 2017 and supplies no named proposer. Its \texttt{corrected\_verified} statement label is supplied metadata, not independent verification here.
\item Supplied report citation: \emph{Low-dimensional Topology and Number Theory} (2018), Oberwolfach Reports, pp.~2363--2426, DOI \texttt{10.4171/owr/2017/38}. \url{https://doi.org/10.4171/owr/2017/38}. The supplied background describes this report as documenting workshop discussions and open problems.
\item Supplied literature reference: \emph{Mahler measure and the Vol Det conjecture}. \url{https://doi.org/10.1112/jlms.12200}.
\item Supplied literature reference: \emph{Low dimensional Topology and Number Theory}. \url{https://doi.org/10.4171/owr/2017/38}.
\item Literature assessment dated exactly \texttt{2026-08-21}: the supplied curation labels the problem open and states that the inequality is proved for several toroidal link examples while the general semi-regular biperiodic alternating case remains open. This is supplied dated curation derived from the OWR material, not independently verified here; its time-sensitive status and cited claims require independent checking.
\end{itemize}
\end{document}
